%\documentclass[11pt,a4paper,twoside]{article}

\usepackage[utf8]{inputenc}

\usepackage[T1,T2A]{fontenc}

\usepackage[english.ancient,russian.ancient]{babel}

\usepackage{graphicx}

\usepackage{array}

\usepackage[doi]{samgtu-bib}

\usepackage{samgtu}

\usepackage{samgtu-name}

\usepackage{mathtools}

\mathtoolsset{showonlyrefs}

\usepackage{desclist}

\usepackage{euscript}

\usepackage{mathrsfs} 

\usepackage{multicol}

\usepackage{mathrsfs}

\usepackage{cite}

\usepackage{cmap}

\usepackage{qrcode}

\allowdisplaybreaks[4]

\usepackage[nodayofweek]{datetime}

\usepackage[thinlines]{easybmat}



\renewcommand{\JournalYear}{2018}

\renewcommand{\JournalVolume}{22}

\renewcommand{\JournalNumber}{4}



\newdate{JournalDateSign}{29}{12}{2018}% Дата подписания Вестника в печать

\renewcommand{\JournalDateSign}{\displaydate{JournalDateSign}}

\newdate{JournalDatePrint}{16}{01}{2019}% Дата выхода Вестника из печати

\renewcommand{\JournalDatePrint}{\displaydate{JournalDatePrint}}

%\renewcommand{\JournalUPL}{16.56} % Усл. печ. л.

%\renewcommand{\JournalUIL}{16.53} % Уч.-изд. л.

\renewcommand{\JournalUPL}{15.24} % Усл. печ. л.

\renewcommand{\JournalUIL}{15.21} % Уч.-изд. л.

\renewcommand{\JournalRegNumber}{220/18} % Регистрационный номер

\renewcommand{\JournalOrderNumber}{2} % Номер заказа



%\renewcommand{\JournalRMonth}{Март} % Месяц выхода Вестника

%\renewcommand{\JournalEMonth}{March} % Месяц выхода Вестника

%\renewcommand{\JournalRMonth}{Июнь} % Месяц выхода Вестника

%\renewcommand{\JournalEMonth}{June} % Месяц выхода Вестника

%\renewcommand{\JournalRMonth}{Сентябрь} % Месяц выхода Вестника

%\renewcommand{\JournalEMonth}{September} % Месяц выхода Вестника

\renewcommand{\JournalRMonth}{Декабрь} % Месяц выхода Вестника

\renewcommand{\JournalEMonth}{December} % Месяц выхода Вестника



\newcommand{\totop}[1]{\raisebox{6pt}[1pt][2pt]{#1}}

\newcommand{\tottop}[1]{\raisebox{12pt}[1pt][2pt]{#1}} 

\newcommand{\totttop}[1]{\raisebox{18pt}[1pt][2pt]{#1}} 

\newcommand{\tottttop}[1]{\raisebox{24pt}[1pt][2pt]{#1}} 

\newcommand{\totttttop}[1]{\raisebox{30pt}[1pt][2pt]{#1}} 

\newcommand{\tobot}[1]{\raisebox{-6pt}[1pt][2pt]{#1}} 

\newcommand{\tobbot}[1]{\raisebox{-12pt}[1pt][2pt]{#1}} 

\newcommand{\tobbbot}[1]{\raisebox{-18pt}[1pt][2pt]{#1}}



\graphicspath{{./00PIC/}}

%\usepackage{samgtu-name}







\renewcommand{\RuCitation}{\noindent\textbf{\CYRO\cyrb\cyrr\cyra\cyrz\cyre\cyrc\ \cyrd\cyrl\cyrya\ \cyrc\cyri\cyrt\cyri\cyrr\cyro\cyrv\cyra\cyrn\cyri\cyrya\\}\selectlanguage{english}{\EnAuthorInContents} {\EnTitleInContents}, \textit{\EnJournalName}\/ [\ParallelJournalName],  \JournalYear, vol.~\JournalVolume, no.~\JournalNumber,  pp.~\thefirstpage--\thelastpage. \doiurl. \selectlanguage{russian}\smallskip}



\renewcommand{\EnCitation}{\noindent\textbf{Please cite this article in press as:\\} {\EnAuthorInContents} {\EnTitleInContents}, \textit{\EnJournalName}\/ [\ParallelJournalName],  \JournalYear, vol.~\JournalVolume, no.~\JournalNumber,  pp.~\thefirstpage--\thelastpage.  \midoiurl. }



\vsgtu{1599} %registration number

\FirstPage{203}

\LastPage{213}







%\pdfcrop

%\draft



\ResearchArticle



%\begin{document}

%

%$\;$

%\vspace{1mm}



%\hfill \qrcode[hyperlink,height=.8in]{http://doi.org/10.14498/vsgtu1599}

%\vspace{-.8in}



\udc{517.984}%

\msc{34A55, 34B24, 47E05}





\rutitle{Об обратной задаче Редже для оператора Штурма–Лиувилля с~отклоняющимся аргументом}



\entitle{On an inverse Regge problem for the Sturm--Liouville operator with deviating argument}







\ruauthor{М.~Ю.~Игнатьев}{И\,г\,н\,а\,т\,ь\,е\,в~М.~Ю.}



\enauthor{M.~Yu.~Ignatiev}{I\,g\,n\,a\,t\,i\,e\,v~M.~Yu.}



\ruaffil{\saratovuni.}



\enaffil{\engsaratovuni.}



\ruauthorsdetails{1}{% Количество авторов

\fullauthor{\rupid{18391}{\selectlanguage{russian}Михаил Юрьевич Игнатьев}}

\CAuthor % Автор-корреспондент

\orcid{0000-0002-4354-9197} % ORCID ID автора 

\regalia{\selectlanguage{russian}кандидат физико-математических наук, доцент} 

\position{\selectlanguage{russian}каф. математической физики и вычислительной математики} 

\email{IgnatievMU@info.sgu.ru}

}





\enauthorsdetails{1}{% Количество авторов

\fullauthor{\enpid{18391}{Mikhail Yu. Ignatiev}}

\CAuthor % Автор-корреспондент

\orcid{0000-0002-4354-9197} % ORCID ID автора 

\regalia{Cand. Phys. \& Math. Sci.; Associate Professor} 

\position{Dept. of Mathematical Physics and Computational Mathematics} 

\email{IgnatievMU@info.sgu.ru}\vspace{-8mm}

}









\rukeywords{дифференциальные операторы, отклоняющийся аргумент, постоянное запаздывание, обратные спектральные задачи, задача Редже}

\enkeywords{differential operators, deviating argument, constant delay, inverse spectral problems,

Regge problem}



\ruabstract{Исследуется краевая задача вида $Ly=\rho^2 y$, $y(0)=y'(\pi)+i\rho y(\pi)=0$, где $L$ "--- оператор Штурма--Лиувилля с постоянным запаздыванием $a$. Данная краевая задача является обобщением классической задачи Редже. Потенциал $q({}\cdot{})$ есть вещественнозначная функция из пространства $L_2(0,\pi)$,  обращающаяся в~0 почти всюду на $(0,a)$. Никаких других ограничений на потенциал не налагается, в частности, не предполагается никаких дополнительных условий относительно поведения $q(x)$ при $x\to \pi$. При столь общих предположениях асимптотическое разложение характеристической функции краевой задачи при $\rho\to\infty$ не содержит главного члена. Как следствие, стандартные методы не позволяют получить в явном виде асимптотику спектра. В работе рассматривается обратная задача восстановления оператора по заданному подмножеству спектра краевой задачи. Обратные задачи для дифференциальных операторов с отклоняющимся аргументом существенно сложнее по сравнению с классическими обратными задачами для дифференциальных операторов. <<Нелокальность>> таких операторов является непреодолимым препятствием для применения классических методов теории обратных задач. Мы рассматриваем обратную задачу в случае запаздывания, большего или равного половине длины интервала, и показываем, что задание подмножества спектра краевой задачи при определенных условиях однозначно определяет потенциал. Соответствующие подмножества спектра описываются в терминах их плотности. В работе также представлена конструктивная процедура решения обратной задачи.}



\enabstract{Boundary value problem of the form $Ly=\rho^2 y$, $y(0)=y'(\pi)+i\rho y(\pi)=0$, where $L$ is the Sturm--Liouville operator with constant delay $a$ is studied. The boundary value problem can be considered as a generalization of the classical Regge problem. The potential $q({}\cdot{})$ is assumed to be a real-valued function from $L_2(0,\pi)$ equal to $0$ a.e. on $(0,a)$. No other restrictions on the potential are imposed, in particular, we make no additional assumptions regarding an asymptotical behavior of $q(x)$ as $x\to\pi$. In this general case, the asymptotical expansion of the characteristic function of the boundary value problem as $\rho\to\infty$ contains no leading term. Therefore, no explicit asymptotics of the spectrum can be obtained using the standard methods.

We consider the inverse problem of recovering the operator from given subspectrum of the boundary value problem.

Inverse problems for differential operators with deviating argument are essentially more difficult with respect to the classical inverse problems for differential operators. ``Non-local'' nature of such operators is insuperable obstacle for classical methods of inverse problem theory.

We consider the inverse problem in case of delay, which is not less than the half length of the interval and establish that the specification of the subspectrum of the boundary value problem determines, under certain conditions, the potential uniquely.

 Corresponding subspectra are characterized in terms of their densities. We also provide a constructive procedure for solving the inverse

problem.}



\newdate{ignatieva}{10}{01}{2018}

\date{\displaydate{ignatieva}}

\newdate{ignatievb}{12}{04}{2018}

\revisiondate{\displaydate{ignatievb}}

\newdate{ignatievc}{11}{06}{2018}

\accepteddate{\displaydate{ignatievc}}



\newdate{ignatieve}{27}{06}{2018}

\renewcommand{\JournalDataOnline}{\displaydate{ignatieve}}



%\ForPeerReview



\vspace{-3mm}



\makeentitle







%%%%%%%%%%%%%%%%%%%%%%%%%%%%%%%%%%%%%%%%%%%%%%%

%Текст статьи

%%%%%%%%%%%%%%%%%%%%%%%%%%%%%%%%%%%%%%%%%%%%%%%



\Section[n]{Introduction}

Consider the boundary value problem:

\begin{gather}

-y''(x)+q(x)y(x-a)=\rho^2 y(x), \quad 0<x<\pi,           

\label{ig:(1)}

\\

y(0)=y'(\pi)+i\rho y(\pi)=0, 

\label{ig:(2)}

\end{gather}

where  $q({}\cdot{})$ is a real-valued function ({\it potential}\/),

$q({}\cdot{})\in L_2(0,\pi)$ and $q(x)=0$ a.e. on $(0,a)$. 

In the paper, we study the inverse problem of recovering the potential from given subspectrum of the boundary value problem \eqref{ig:(1)}, \eqref{ig:(2)}. 

Inverse spectral problems consisting in recovering operators from their spectral characteristics often appear in mathematics, physics, mechanics, geophysics, electronics, meteorology, etc. The greatest success in the inverse

problem theory has been achieved for the classical Sturm–Liouville operator and afterwards for higher order differential operators (see \cite{ign:1, ign:2, ign:3, ign:4}  and references therein). At the same time, the classical methods of inverse spectral theory, which allow obtaining global solutions

of inverse problems for differential operators, are not applicable for differential operators

with deviating argument as well as for other classes of nonlocal operators such as integro-differential,

integral and other operators \cite{ign:5, ign:6, ign:7}. Therefore, the general inverse spectral theory for

differential operators

with deviating argument has not yet been constructed and there are only isolated results in this

direction not forming the general picture \cite{ign:8, ign:9, ign:10, ign:11} (see also references therein).



\smallskip



\Section{Preliminary Information}

We note first that the spectrum of the boundary value problem \eqref{ig:(1)}, \eqref{ig:(2)} coincides with the set of zeros of the characteristic function  $\Delta(\rho)=S'(\pi,\rho)+i\rho S(\pi,\rho)$, where $S(x,\rho)$ is a solution of \eqref{ig:(1)} under the initial conditions $S(0,\rho)=0$, $S'(0,\rho)=1$. 

If $q(x)=0$ a.e. on $(0,\pi)$, then $\Delta(\rho)=\Delta_0(\rho)=\exp(i\rho \pi)$ and the corresponding spectrum is empty. 

Converse proposition is also true.



\smallskip



\hypertarget{ign:th1}{}



{\small \sc Theorem~1.} {\it If the spectrum of the boundary value problem {\rm \eqref{ig:(1)}, \eqref{ig:(2)}} is empty$,$ then $q(x)=0$ a$.$e$.$ on $(0,\pi).$}



\smallskip



The proof of the Theorem \hyperlink{ign:th1}{1} is actually contained in the proof of the uniqueness theorem in \cite{ign:5}.





\smallskip



Everywhere below we assume that $a\in[\pi/2,\pi)$. Then for $x\ge a$ the following representations hold:

\begin{gather}

S(x,\rho)=\frac{\sin\rho x}{\rho}-\frac{\cos\rho(x-a)}{2\rho^2}

\int_a^x q(t)\,dt+\frac{1}{2\rho^2}\int_a^x q(t)\cos\rho(x-2t+a)\,dt,

\\

S'(x,\rho)=\cos\rho x+\frac{\sin\rho(x-a)}{2\rho} \int_a^x q(t)\,dt

- \frac{1}{2\rho}\int_a^x q(t)\sin\rho(x-2t+a)\,dt.

\end{gather}

Therefore, we have:

\begin{multline}

\Delta(\rho)=\exp(i\rho\pi)

+\frac{1}{2i\rho}\exp(i\rho(\pi-a))\int_a^\pi q(t)\,dt -\\

\frac{1}{2i\rho}\exp(i\rho(\pi+a))\int_a^\pi q(t)\exp(-2i\rho t)\,dt.

\end{multline}



As in the case of the classical Regge problem (see, for instance, \cite{ign:12} and references therein) the asymptotical expansion of the characteristic function for

 $\rho\to\infty$, $\mbox{Im}\rho>0$ does not contain a leading term and the standard methods do not allow obtaining the explicit asymptotic formulas for the spectrum without additional restrictions on $q({}\cdot{})$.





In this paper, we study the inverse problem in a general case, without any additional restrictions on $q({}\cdot{})$. In this case the properties of the spectrum can be characterized in terms of its angular density.

For a set $M$ we denote by $n_M(t)$ the number of elements of $M$ in the circle $\{|z|<t\}$. We recall that the first order countable set $M$ is referred to have an angular density if for all pairs $\theta_1<\theta_2$ (with possible exception of some countable set) there exists the limit:

$$

P(\theta_1,\theta_2)=\lim\limits_{t\to\infty}\frac{n(t;\theta_1,\theta_2)}{t},

$$

where $n(t;\theta_1,\theta_2)$ is the number of elements of $M$ in the sector $\bigl\{z=r\exp(i\theta): r\in(0,t),$ $ \theta\in(\theta_1,\theta_2)\bigr\}$. 

We say that the angular density of the set $M$ is determined by the nondecreasing function $D(\theta)$ if $P(\theta_1,\theta_2)=D(\theta_2)-D(\theta_1)$. Suppose the first order countable set $M$ has the angular density. We call $M$ a properly distributed set if the series  $\sum_{z\in M}(1/z)$ converges in the principal-value sense, i.e., the following limit exists:

$$

\lim\limits_{r\to\infty}\sum\limits_{z\in M, |z|< r}\frac{1}{z}.

$$



Our analysis is based on the connection between the properties of the roots of entire functions and their growth. For reader's convenience we recall here some of theorems that we shall use below (see \cite{ign:13, ign:14}  for details).



\smallskip



\hypertarget{ign:th2}{}



{\small \sc Theorem~2 (Cartwright theorem).} {\it Let $\{z_n\}$ be the set of all nonzero roots of the entire function

of exponential type $F(z).$ If $F(z)$ is bounded on the real axis$,$ then${:}$

 \begin{enumerate}

   \item[\rm 1)] $F({}\cdot{})$ have the complete regularity of growth and its  indicator is given by the formula{\rm:}

   $$h(\theta)=h_\pm |\sin\theta|, \quad \pm\theta\in[0,\pi],$$

   where $h_+,$ $ h_-$ are real values such that $h_+ +h_-\geq 0,$ if$,$ in addition$,$ the set $\{z_n\}$ is nonempty$,$ then $h_+ +h_->0;$

   \item[\rm 2)] $\sum_n \left|\mbox{\rm Im}\frac{1}{z_n}\right|<\infty;$

   \item[\rm 3)] for any $\varepsilon\in(0,\pi/2)$ the angular density of the first order sequence $\{z_n\}$ in each of angles $|\arg z|<\varepsilon,$ $|\pi-\arg z|<\varepsilon$ is equal to $(h_+ +h_-)/2\pi,$ the sequence $\{z_n\}$ has at most finite number of terms outside these angles$;$

   \item[\rm 4)] the series $\sum_n \frac{1}{z_n}$ converges in the principal-value sense$.$

 \end{enumerate}

 }



\smallskip



\hypertarget{ign:th3}{}



{\small \sc Theorem~3.} {\it

 Suppose the set of roots $\{z_n\}$ of the canonical product

 $$

 f(z)=\prod\limits_n \Bigl(1-\frac{z}{z_n}\Bigr)\exp\Bigl(\frac{z}{z_n}\Bigr)

 $$

 is a properly distributed set with the angular density determined by the function $D({}\cdot{}).$ 

 Then for any $z$ outside some $C^0$-set the following asymptotics holds

 $$

 \ln\bigl|f\bigl(r\exp(i\theta)\bigr)\bigr|=H(\theta)\bigl(r+o(r)\bigr), \quad r\to\infty,

 $$

 where

 $$

 H(\theta)=\int_{\theta-2\pi}^\theta (\theta-\psi)\sin(\theta-\psi-\pi)dD(\psi)+\tau_0\cos(\theta-\theta_0),

 $$

 $$

 \tau_0\exp(i\theta_0)=\delta_0:=\sum\limits_{n}\frac{1}{z_n}.

 $$

 }

 

 \smallskip



Consider again the boundary value problem \eqref{ig:(1)}, \eqref{ig:(2)}.

One can write the representation for its characteristic function as follows:

\begin{gather}

\label{ig:(3)}

\Delta(\rho)=\exp(i\rho\pi)+\int_{-(\pi-a)}^{\pi-a} u(t) \exp(i\rho t)\,dt , 

\\

u(t)=\frac{1}{2}\int_{({\pi+a-t})/{2}}^\pi q(\xi)\, d\xi. 

\label{ig:(4)}

\end{gather}



From representation \eqref{ig:(3)}, \eqref{ig:(4)} and the Paley--Wiener theorem  it follows that

   $\Delta({}\cdot{})$ is an entire function of exponential type, it has a complete regular growth and its  conjugate diagram is the segment

 $[i(\pi+a-2b), i\pi]$, where $b:=\sup  {\rm supp}\ q$.

Moreover, if  $q(x)$ is real-valued, then $\overline{\Delta(-\bar \rho)}=\Delta(\rho)$.



Using the Cartwright theorem we arrive at the following result.



\smallskip



\hypertarget{ign:proposition1}{}



{\small \sc Proposition~1.} {\it If $q({}\cdot{})$ is not equal to $0$ identically then the problem \eqref{ig:(1)}, \eqref{ig:(2)} has a countable set $\sigma$ of eigenvalues$.$ The set $\sigma$ is a properly distributed set$,$ its angular density is  determined by the function 

$$

D(\theta)=\frac{2b-a}{2\pi}\Bigl[\frac{\theta}{\pi}\Bigr] ,

$$

   where $[{}\cdot{}]$ denote the integer part$.$ 

   If $q(x)$ is real-valued$,$ then the set $\sigma$ is symmetric with respect to the imaginary axis$,$ moreover$,$ the set of pure imaginary eigenvalues is $($at most\/$)$ finite$.$

 }



\smallskip



{\small\sc Remark~1.} In formulations of Theorems \hyperlink{ign:th2}{2}, \hyperlink{ign:th3}{3} the roots $\{z_n\}$ are counted according to their algebraic multiplicities. Similarly, the terms in the spectrum $\sigma$ are counted according to their algebraic multiplicities.



\smallskip











\Section{Inverse Problem}

For the classical inverse Regge problem, it is known that the spectrum of the boundary value problem \eqref{ig:(1)}, \eqref{ig:(2)}  (with $a=0$) uniquely determines the potential. 

In the case of $a\geq\pi/2$ it will be found that the specification of some {\it part} of the spectrum uniquely determines the potential. 

We characterize the proper subspectra in terms of their densities. 

Such a method of characterization of the subspectra was introduced in \cite{ign:15}, where  the so called {\it incomplete} inverse problems were considered.



\smallskip



\hypertarget{ign:problem1}{}



{\small \sc Problem 1.} {\it Given the symmetric with respect to the imaginary axis subspectrum $\Lambda\subset\sigma$ of the boundary value problem {\rm \eqref{ig:(1)}, \eqref{ig:(2)}, } find the potential $q(x),$ $ x\in[a,\pi].$}





\smallskip



The set $\Lambda$ possibly contains multiple elements but the multiplicity of each element $\rho_0\in\Lambda$ is less or equal to the algebraic multiplicity of $\rho_0$ as an eigenvalue of the boundary value problem \eqref{ig:(1)}, \eqref{ig:(2)}. For the definiteness we assume $0\notin\Lambda$.



We establish first the uniqueness result for the solution of the Problem~\hyperlink{ign:problem1}{1}.

We agree that if a certain symbol $\xi$ denotes an object related to the boundary value problem \eqref{ig:(1)}, \eqref{ig:(2)} with the potential $q({}\cdot{})$, then $\tilde\xi$ will denote the analogous object related to the boundary value problem \eqref{ig:(1)}, \eqref{ig:(2)}  with the potential $\tilde q({}\cdot{})$, and $\hat\xi:=\xi-\tilde\xi$.





\smallskip



\hypertarget{ign:th4}{}



{\small \sc Theorem~4.} {\it Suppose there exists $($and is given\/$)$ the limit $$ \lim\limits_{r\to\infty} \frac{n_{\Lambda}(r)}{n_{\sigma}(r)}=K.$$

Suppose that $$K>\frac{2\pi-2a}{2\pi-a}.$$ Then the specification of the subspectrum $\Lambda$ determines the potential uniquely$.$ Namely$,$ if $\Lambda\subset\sigma,$ $\Lambda\subset\tilde\sigma$ and

$$

\lim\limits_{r\to\infty} \frac{n_{\Lambda}(r)}{n_{\sigma}(r)}=\lim\limits_{r\to\infty} \frac{n_{\Lambda}(r)}{n_{\tilde\sigma}(r)}=

K>\frac{2\pi-2a}{2\pi-a},

$$

then $q(x)=\tilde q(x)$ a$.$e$.$ on $[a,\pi].$

}



\smallskip





\textit{Proof.}

We note first that the specification of $K$ determines uniquely the density of the spectrum $\sigma$ via the relation:

$$

\lim\limits_{r\to\infty}\frac{n_\sigma(r)}{r}=K^{-1}\cdot\lim\limits_{r\to\infty}\frac{n_\Lambda(r)}{r}.

$$

By virtue of Proposition~\hyperlink{ign:proposition1}{1} this means that the specification of $K$ determines uniquely the value of $b$ that can be calculated by the formulas:

\begin{equation}

b=\frac{a+\pi K_1}{2}, \quad  K_1=K^{-1}\cdot\lim\limits_{r\to\infty}\frac{n_\Lambda(r)}{r}. 

\label{ig:(5)}

\end{equation}

In what follows we assume $b$ to be known.



 Suppose $\Lambda=\{\rho_n\}_{n\geq 1}$. Without loss of generality we assume that multiple elements of $\Lambda$ are neighboring. We define the functions $e_n(x)$ as follows. If $\rho_n\neq \rho_{n-1}$ (or $n=1$), then $e_n(x):=\exp(i\rho_n x)$. 

 If $\rho_n=\rho_{n-1}=\dots=\rho_{n-\nu}\neq\rho_{n-\nu-1}$, then

$$

e_n(x):=\frac{d^\nu}{d\rho^\nu}\bigl(\exp(i\rho x)\bigr)\Bigr|_{\rho=\rho_n}.

$$

Let $\sigma$, $\tilde\sigma$ be the spectra of the boundary value problems \eqref{ig:(1)}, \eqref{ig:(2)}  with the potentials $q$ and $\tilde q$ respectively. 

We assume (see the discussion above) that $\sup  {\rm supp}\ q=\sup {\rm supp}\ \tilde q =b$. 

By virtue of \eqref{ig:(3)}, \eqref{ig:(4)}   $\Lambda\subset\sigma$ imply the relations:

$$

\frac{d^\nu}{d\rho^\nu}

\biggl( \exp(i\rho \pi)+\exp\bigl(i\rho(\pi-b)\bigr)

\int_{-(b-a)}^{b-a} v(t)\exp(i\rho t)dt\biggr)\biggr|_{\rho=\rho_n}\!\!\!\! =0, \; \; n=1,2,\dots,

$$

where

$$

v(t)=u(t+\pi-b)

$$

and, as above, for each $n$ $\nu$ is such that $\rho_n=\rho_{n-1}=\dots=\rho_{n-\nu}\neq\rho_{n-\nu-1}$.

Similarly, we have:

$$

\frac{d^\nu}{d\rho^\nu}

\biggl(

\exp(i\rho \pi)+\exp\bigl(i\rho(\pi-b)\bigr)

\int_{-(b-a)}^{b-a} \tilde v(t)\exp(i\rho t)dt\biggr)\biggr|_{\rho=\rho_n}\!\!\!\! =0, \; \; n=1,2,\dots.

$$

Therefore, one can write:

\begin{equation}

\label{ig:(6)}

\int_{-(b-a)}^{b-a} \hat v(t)e_n(t)dt=0, \quad  n=1,2,\dots. 

\end{equation}

Let us show that the system $\{e_n(\cdot)\}_{n\geq 1}$ is complete in $L_2(a-b,b-a)$. 

By virtue of the Levinson theorem   \cite{ign:13, ign:14}, it is sufficient to prove that

\begin{equation}

\label{ig:(7)}

l(\Lambda):=\overline{\lim\limits_{R\to\infty}}\Bigl(N_\Lambda(R)-\frac{2(b-a)}{\pi}R+\frac{1}{2}\log R\Bigr)>-\infty. 

\end{equation}

Here

$$

N_\Lambda(R)=\int_0^R\frac{n_\Lambda(r)}{r}\,dr.

$$

Given arbitrary  $\varepsilon>0$, for $r>r_0(\varepsilon)$ we have:

$$

n_\sigma(r)>(1-\varepsilon)\frac{2b-a}{\pi}r,

$$

$$

n_\Lambda(r)>(1-\varepsilon)Kn_\sigma(r)>(1-\varepsilon)^2(1+\beta)\frac{2\pi-2a}{2\pi-a}\frac{2b-a}{\pi} r,

$$

where $\beta>0$. Therefore,

$$

\frac{n_\Lambda(r)}{r}>(1-\varepsilon)^2(1+\beta) \frac{2b-2a}{\pi}  \frac{2\pi-2a}{2\pi-a} \frac{2b-a}{2b-2a}=:A.

$$

Under the conditions of the theorem, by choosing $\varepsilon{>}0$ such that $(1-\varepsilon)^2{(1+\beta){>}1} $  and taking into account that

$$

\frac{2\pi-2a}{2\pi-a} \frac{2b-a}{2b-2a}\geq 1

$$

 we obtain $$A>\frac{2b-2a}{\pi}.$$ 

 Then for all sufficiently large $R>0$ one has:

$$

N_\Lambda(R)\geq A(R-r_0)>\frac{2(b-a)}{\pi}R

$$

and condition \eqref{ig:(7)} is fulfilled.



Since the system $\{e_n(\cdot)\}_{n\geq 1}$ is complete, from \eqref{ig:(6)}  it follows that $\hat v(x)=0$ a.e. on $(a-b,b-a)$, therefore $u(x)=\tilde u(x)$, $q(x)=\tilde q(x)$.

\hfill$\Box$



\smallskip



Our next goal is to provide a constructive procedure for solving  Problem~\hyperlink{ign:problem1}{1}. 

Suppose that the subspectrum $\Lambda=\{\rho_n\}_{n\geq 1}$ is given and the conditions of Theorem 4 are satisfied. This means, in particular, that we can assume the value of $b=\sup {\rm supp}\ q$ to be also given.



We choose $\tilde q=0$. Then:

\begin{gather}

\hat\Delta(\rho)=\int_{-(\pi-a)}^{\pi-a} u(t) \exp(i\rho t)\,dt=\exp\bigl(i\rho(\pi-b)\bigr)\int_{-(b-a)}^{b-a} v(t)\exp(i\rho t)\,dt, 

\label{ig:(8)}

\\

v(t)=u(t+\pi-b).

\label{ig:(9)}

\end{gather}



Let us note that under the conditions of Theorem~\hyperlink{ign:th4}{4}

$\Lambda$ is a properly distributed set with the angular density determined by the function:

$$

D(\theta)=K\frac{2b-a}{2\pi}\Bigl[\frac{\theta}{\pi}\Bigr].

$$

Indeed, to show this we actually need only to prove the existence of the limit:

$$

\lim\limits_{r\to\infty}\sum\limits_{|\rho_n|< r} \frac{1}{\rho_n}.

$$

Since $\Lambda$ is symmetric with respect to the imaginary axis, we have:

$$

\sum\limits_{|\rho_n|< r} \mbox{Re}\frac{1}{\rho_n}=0,

$$

therefore there exists the limit:

$$

\lim\limits_{r\to\infty}\sum\limits_{|\rho_n|< r} \mbox{Re}\frac{1}{\rho_n}=0.

$$

Furthermore, by virtue of the Cartwright theorem one has:

$$

\sum\limits_{\rho\in\sigma}\Bigl|\mbox{Im}\frac{1}{\rho}\Bigr|<\infty,

$$

and therefore the series

$$

\sum\limits_{n=1}^\infty \mbox{Im}\frac{1}{\rho_n}

$$

converges absolutely.



Define the function:

\begin{equation}

\omega(\rho):=\exp(i\kappa\rho)\prod\limits_{n=1}^\infty

\Bigl(1-\frac{\rho}{\rho_n}\Bigr)\exp\Bigl(\frac{\rho}{\rho_n}\Bigr), 

\label{ig:(10)}

\end{equation}

 where the real number $\kappa$ will be chosen below. 

 By virtue of Theorem~\hyperlink{ign:th3}{3} under the conditions of Theorem~\hyperlink{ign:th4}{4} $\omega({}\cdot{})$ is an entire function of exponential type and its indicator is:

$$

H_\omega(\theta)=(b-a/2)K|\sin\theta|+(\tau_0-\kappa)\sin\theta,

$$

where

$$

\tau_0=\sum\limits_{n=1}^\infty \mbox{Im}\frac{1}{\rho_n}.

$$

Now let us choose $\kappa$ such that the conjugate diagram of the function $\omega({}\cdot{})$ will strictly contain the segment $[i(\pi+a-2b), i(\pi-a)]$, 

that is, by virtue of \eqref{ig:(8)}, the conjugate diagram of the function $\hat\Delta({}\cdot{})$. 

Note that such a value of $\kappa$ exists since

$$

H_\omega(-\pi/2)+H_\omega(\pi/2)=(2b-a)K> 2b-2a=H_{\hat\Delta}(-\pi/2)+H_{\hat\Delta}(\pi/2)

$$

for any $\kappa$. 

By virtue of the Paley--Wienner theorem and Theorems~\hyperlink{ign:th2}{2},~\hyperlink{ign:th3}{3} one has $H_\omega(\theta)>H_{\hat\Delta}(\theta)$ for all $\theta$. 

Moreover, by virtue of Theorem~\hyperlink{ign:th3}{3}, for $\rho\to\infty$ outside some $C^0$-set we have:

\begin{equation}

\log|\omega(\rho)|=H_\omega(\theta)\bigl(r+o(r)\bigr), \quad  \rho=r\exp(i\theta). 

\label{ig:(11)}

\end{equation}

Since all the circles that do not contain the points of $\Lambda$ can be excluded from the $C^0$-set \cite{ign:13}, asymptotics \eqref{ig:(11)} holds for $\rho\to\infty$, $\rho\in G_\delta:=\{\rho:\mbox{dist}(\rho,\Lambda)>\delta\}$ with arbitrary $\delta>0$. 

For all sufficiently small $\delta>0$ one can choose the subsequence of the circles $\{|\rho|=R_N\}$, $R_N\to\infty$ lying in $G_\delta$. 

Our construction yields:

\begin{equation}

\lim\limits_{N\to\infty}\max\limits_{|\rho|=R_N}\Bigl|\frac{\hat\Delta(\rho)}{\omega(\rho)}\Bigr|=0. 

\label{ig:(12)}

\end{equation}

From \eqref{ig:(12)} it follows that:

$$

\lim\limits_{N\to\infty}\int\limits_{|\mu|=R_N}\frac{\hat\Delta(\mu)}{\omega(\mu)}\frac{d\mu}{\rho-\mu}=0,

$$

and the residue theorem yields:

\begin{equation}

\frac{\hat\Delta(\rho)}{\omega(\rho)}=

\sum\limits_{\mu_0\in\Lambda_0} \mathop{\rm Res}_{\mu=\mu_0}

\Bigl(\frac{\hat\Delta(\mu)}{\omega(\mu)(\rho-\mu)}\Bigr):=

\lim\limits_{N\to\infty}\sum\limits_{\mu_0\in\Lambda_0(N)} \mathop{\rm Res}_{\mu=\mu_0}\Bigl(\frac{\hat\Delta(\mu)}{\omega(\mu)(\rho-\mu)}\Bigr), 

\label{ig:(13)}

\end{equation}

where

$\Lambda_0$ is the set with no multiple elements containing all the (different) elements of $\Lambda$,

$\Lambda_0(N):=\Lambda_0\cap\{|\mu|<R_N\}$.



Define the functions:

\begin{equation}

\omega_{n}(\rho)=\frac{1}{2\pi i}\int_{|\mu-\rho_n|=\delta}\frac{(\mu-\rho_n)^\nu}{(\rho-\mu)\omega(\mu)}\,d\mu, \quad  \rho\in\mathbb{C}\setminus\Lambda_0,

\label{ig:(14)}

\end{equation}

 where $\nu$ is such that $\rho_n=\dots=\rho_{n-\nu}\neq\rho_{n-\nu-1}$ ($\nu=0$ if $\rho_n\neq\rho_{n-1}$ or $n=1$) and $\delta\in(0,|\rho-\rho_n|)$ is such that the circle $\{\mu:|\mu-\rho_n|\leq\delta\}$ does not contain the points of $\Lambda_0$ different from $\rho_n$. 

 Then for any $\mu_0=\rho_n=\dots=\rho_{n+m_n-1}\neq\rho_{n+m_n}$, $\mu_0\neq\rho_{n-1}$ we have:

\begin{equation}

\mathop{\rm Res}_{\mu=\mu_0}\Bigl(\frac{\hat\Delta(\mu)}{\omega(\mu)(\rho-\mu)}\Bigr)=\sum\limits_{\nu=0}^{m_n-1}

\frac{\hat\Delta^{(\nu)}(\rho_n)}{\nu !}\omega_{n+\nu}(\rho). 

\label{ig:(15)}

\end{equation}

Furthermore, since $\rho_n$ is an eigenvalue of the boundary value problem \eqref{ig:(1)}, \eqref{ig:(2)}  and its algebraic multiplicity is not less than  $m_n$ we have $\Delta^{(\nu)}(\rho_n)=0$, $\nu=\overline{0,m_n-1}$. 

By virtue of \eqref{ig:(3)}  this yields:

\begin{equation}

\frac{\hat\Delta^{(\nu)}(\rho_n)}{\nu !}=-\frac{1}{\nu !} \frac{d^{\nu}}{d\rho^\nu}\bigl(\exp(i\rho\pi)\bigr)\Bigr|_{\rho=\rho_n}=:\beta_{n+\nu}. 

\label{ig:(16)}

\end{equation}

Then \eqref{ig:(15)} can be written as:

$$

\mathop{\rm Res}_{\mu=\mu_0}\Bigl(\frac{\hat\Delta(\mu)}{\omega(\mu)(\rho-\mu)}\Bigr)

=\sum\limits_{\nu=0}^{m_n-1}

\beta_{n+\nu}\omega_{n+\nu}(\rho).

$$

Substituting this into \eqref{ig:(13)} we obtain:

$$

\frac{\hat\Delta(\rho)}{\omega(\rho)}=\sum\limits_{n=1}^\infty \beta_n \omega_n(\rho),

$$

where the series in the right-hand side converges in the principal-value sense: $$\sum_{n=1}^\infty:=\lim_{N\to\infty}\sum_{|\rho_n|<R_N}.$$ 

Finally, by multiplying both sides of the  relation by $\omega(\rho)$, we obtain the following reconstruction formula for the function $\hat\Delta({}\cdot{})$:

\begin{equation}

\hat\Delta(\rho)=\omega(\rho)\sum\limits_{n=1}^\infty \beta_n \omega_n(\rho).

\label{ig:(17)}

\end{equation}

Now we use representations \eqref{ig:(8)}, \eqref{ig:(9)}  to complete our procedure of solving Problem~\hyperlink{ign:problem1}{1}. Thus, we arrive at the following result.



\smallskip



{\small\sc Theorem~5.} {\it Under the conditions of Theorem~{\rm \hyperlink{ign:th4}{4}} the potential $q(x),$ $x\in[a,\pi]$ can be recovered by making the following steps\/${:}$

\begin{enumerate}

  \item[\rm 1.] Find $b$ by formula {\rm \eqref{ig:(5)}.}

  \item[\rm 2.] Construct the function $\omega({}\cdot{})$ by formula {\rm \eqref{ig:(10)}.}

  \item[\rm 3.] Calculate $\{\beta_n\}$ and construct the functions $\{\omega_n({}\cdot{})\}$ using \eqref{ig:(16)} and \eqref{ig:(14)}  respectively$.$

  \item[\rm 4.] Recover the function $\hat\Delta({}\cdot{})$ via {\rm  \eqref{ig:(17)}.}

  \item[\rm 5.] Find the $u({}\cdot{})$ as the inverse Fourier transform of the function $\hat\Delta({}\cdot{}).$ Recover $q({}\cdot{})$ by the formula   $q(x)=4u'(\pi+a-2x).$

\end{enumerate}



}



\AdditionalContent{Competing interests}{I declare that I have no competing interests.} 

\AdditionalContent{Author's Responsibilities}{I take full responsibility for submitting the final manuscript in print. I approved the final version of the manuscript.} 

\AdditionalContent{Funding}{This work was supported by the Russian Foundation for

Basic Research (projects nos.~16--01--00015\_a, 17--51--53180\_GPhEN\_a) and Russian Ministry of Edu\-cation and Science (project no.~1.1660.2017/4.6).}















\begin{thebibliography}{99}



\Bibitem{ign:1}

\by Marchenko~V.~A.

\book Sturm--Liouville operators and applications

\zmath{0592.34011}

\serial Operator Theory: Advances and Applications

\vol 22

\publaddr Basel, Boston, Stuttgart

\publ Birkh\"auser Verlag

\totalpages xi+367 

\yr 1986

\moreref Russ. ed.:

\publaddr Kiev

\publ Naukova Dumka

\yr 1977

\totalpages 393 

\zmath{0399.34022}





\Bibitem{ign:2}

\by Levitan~B.~M.

\book Inverse Sturm--Liouville problems

\zmath{0749.34001}

\publaddr Utrecht

\publ VNU Science Press

\totalpages x+240 

\yr 1987

\moreref Russ. ed.:

\zmath{0575.34001}

\publaddr Moscow

\publ Nauka

\yr 1984

\totalpages 246 



\Bibitem{ign:3}

\by Beals~R.,  Deift~P., Tomei~C.

\book Direct and inverse scattering on the line

\zmath{0679.34018}

\serial Mathematical Surveys and Monographs

\vol 28

\publaddr Providence, RI

\publ American Mathematical Society 

\totalpages xiii+209 

\yr 1988





\Bibitem{ign:4}

\by Yurko~V.~A.

\book Introduction to the theory of universe spectral problems

\lang In Russian

\zmath{1137.34001}

\publaddr Moscow

\publ Fizmatlit 

\totalpages 384 

\yr 2007



\Bibitem{ign:5}

\by Buterin~S.~A.

\jour Result. Math. 

\yr 2007

\vol 50

\crossref{10.1007/s00025-007-0244-6}

\paper On an inverse spectral problem for a convolution integro-differential operator

\issue 3--4

\pages 73–181



\Bibitem{ign:6}

\by Kuryshova~Ju.~V. 

\paper  Inverse spectral problem for integro-differential operators

\jour Math. Notes

\crossref{10.1134/S0001434607050240}

\yr 2007

\vol 81

\issue 6

\pages 767–777



\Bibitem{ign:7}

\by Bondarenko~N.~P., Buterin~S.~A.

\paper On recovering the Dirac operator with an integral delay from the spectrum

\jour Result. Math. 

\crossref{10.1007/s00025-016-0568-1}

\yr 2017

\vol 71

\issue 3

\pages 1521–1529



\Bibitem{ign:8}

\by Freiling~G., Yurko~V.~A.

\paper Inverse problems for Sturm--Liouville differential operators with a constant delay 

\jour Appl. Math. Lett. 

\yr 2012

\vol 25

\crossref{10.1016/j.aml.2012.03.026}

\issue 11

\pages 1999–2004





\Bibitem{ign:9}

\by Vladi\v{c}i\'{c}~V., Pikula~M. 

\paper An inverse problem for Sturm–Liouville-type differential equation with a

constant delay

\crossref{10.5644/SJM.12.1.06}

\jour Sarajevo J. Math.

\yr 2016

\vol 12(24)

\issue 1

\pages 83–88





\Bibitem{ign:10}

\by Buterin~S.~A., Pikula~M., Yurko~V.~A. 

\paper Sturm--Liouville differential operators with deviating argument

\crossref{10.5556/j.tkjm.48.2017.2264}

\jour Tamkang J. Math.

\yr 2017 

\vol 48

\issue 1

\pages 61–71



\Bibitem{ign:11}

\by Yurko~V.~A., Buterin~S.~A. 

\paper An inverse spectral problem for Sturm–Liouville operators with a large constant delay

\jour Anal. Math. Phys. 

\yr 2017

\crossref{10.1007/s13324-017-0176-6}



\Bibitem{ign:12}

\by Gubreev~G.~M., Pivovarchik~V.~N.

\paper Spectral analysis of the Regge problem with parameters

\crossref{10.1007/BF02466004}

\jour Funct. Anal. Appl.

\yr 1997

\vol 31

\issue 1 

\pages 54–57





\Bibitem{ign:13}

\by Levin~B.~Ya.

\zmath{https://zbmath.org/?q=an:0152.06703}

\book Distribution of zeros of entire functions

\serial Translations of Mathematical Monographs

\publ Amer. Math. Soc. 

\publ Providence, R.I.

\yr 1964

\totalpages viii+493 

\moreref Russ. ed.:

\publaddr Moscow

\publ Gostechizdat

\yr 1956

\totalpages 632

\zmath{0111.07401}



\Bibitem{ign:14}

\by Sedletskiy~A.~M. 

\book  Classes of analytic Fourier transforms and exponential approximations

\lang In Russian

\publaddr Moscow

\publ Fizmatlit

\yr 2005

\totalpages 504



\Bibitem{ign:15}

\by Gesztesy~F., Simon~B. 

\paper Inverse spectral analysis with partial information on the potential, II. The case of discrete spectrum

\jour Trans. Amer. Math. Soc. 

\crossref{10.1090/S0002-9947-99-02544-1}

\vol 352 

\issue 6

\yr 2000

\pages 2765--2787 



\end{thebibliography}



\renewcommand{\baselinestretch}{0.89}



\selectlanguage{russian}



\vfill



\AdditionalContent{Конкурирующие интересы}{Конкурирующих интересов не имею.}



\AdditionalContent{Авторский вклад и ответственность}{Я несу полную ответственность за предоставление окончательной версии рукописи в печать. Окончательная версия 

рукописи мною одобрена.}



\AdditionalContent{Финансирование}{Работа выполнена при частичной финансовой поддержке \mbox{РФФИ}

(проекты № 16--01--00015\_a, 17--51--53180\_ГФЕН\_а) и Министерства образования и~науки Российской Федерации (проект 1.1660.2017/4.6).}







\newpage







$\;$

\vspace{-7mm}







\makerutitle









\renewcommand{\baselinestretch}{0.9}



\renewcommand{\RuCitation}{\noindent\textbf{\CYRO\cyrb\cyrr\cyra\cyrz\cyre\cyrc\ \cyrd\cyrl\cyrya\ \cyrc\cyri\cyrt\cyri\cyrr\cyro\cyrv\cyra\cyrn\cyri\cyrya\\} {\RuAuthorInContents} {\RuTitleInContents}~/\!/ \textit{\RuJournalName}, \JournalYear.  \CYRT.~\JournalVolume, \textnumero~\JournalNumber.  \CYRS.~\thefirstpage--\thelastpage. \doiurl.}



\renewcommand{\EnCitation}{\noindent\textbf{Please cite this article in press as:\\} {\EnAuthorInContents} {\EnTitleInContents}, \textit{\EnJournalName}\/ [\ParallelJournalName],  \JournalYear, vol.~\JournalVolume, no.~\JournalNumber,  pp.~\thefirstpage--\thelastpage.  \midoiurl\ (In Russian).}









%\end{document}

